\section{Discussion}
\label{sec:discussion}

\subsection{Implications}

The results indicate that heuristic static analysis can produce useful specifications for 99.0\% of methods in the evaluation corpus. Such specifications support machine-checkable documentation~\cite{peters1998using}, onboarding aids for understanding method contracts~\cite{sillito2008questions}, and signals for tracking API evolution~\cite{tip2002practical}. With an average analysis time of 88\,ms per file, integration into continuous-integration pipelines (pre-commit hooks, pull-request analysis, automated test generation) is practicable.

The comparison with Jdoctor reveals complementary strengths: Jdoctor recovers documented exceptional behaviour from Javadoc, whereas JML-Inferrer derives specifications from code structure and is therefore effective on undocumented methods. A practical workflow may combine both to maximise coverage.

Specifications improve test generation along three independent dimensions. Preconditions delimit valid input ranges, guiding the construction of boundary tests; postconditions define expected outputs, enabling oracles stronger than ``does not throw''; and the clause-by-clause decomposition allows LLMs to target individual specification clauses rather than to reason about entire implementations. Although source-code access (P4) yields higher absolute scores, the substantial improvement from P1 to P3 confirms that specifications encode behavioural information that signatures alone cannot convey.

\paragraph{Uses beyond test generation.} The same inferred specifications serve at least three downstream uses outside the LLM test-generation setting that has been the primary focus of this article. First, they act as the callee contracts that compositional verification of caller code requires: a modular verifier reasoning about a method that calls \texttt{m} requires only \texttt{m}'s pre- and postcondition, not its body, and the inferred contracts supply that summary at library scale without manual annotation effort. A change to a library entry implies a re-check of its callers rather than a re-analysis of the entire transitive dependency graph. Second, they function as a machine-readable contract that AI coding agents can consume when interacting with libraries that postdate the model's training cut-off: an agent asked to use, say, a library released after its training window has no internal model of the library's contracts, but it can read the JML annotations directly and reason against them. Third, they strengthen the oracles of any test-generation tool that consumes JML, not only LLM-driven ones---tools such as JMLUnit and specification-aware EvoSuite have always required contracts as input, and the inferrer supplies the missing upstream step. The empirical evaluation in this article isolates LLM test generation because mutation testing makes that effect measurable, but the same artefacts are the input to the other two uses.

These implications generalise beyond the Apache Commons Lang corpus in principle. The technique is applicable to any Java codebase, with effectiveness depending on code structure: well-structured code yields stronger specifications, and domain-specific business logic may require extending the heuristic repertoire. The JML specification format is independent of any particular LLM and widely recognised~\cite{leavens2006design}; comparable improvements are therefore expected with other models. Application areas extend to library API specification, legacy code modernisation, documentation generation, and runtime verification. Empirical confirmation across heterogeneous codebases remains future work and is identified explicitly as an external-validity threat in Section~\ref{sec:threats}.

\subsection{Limitations}

\paragraph{Specification coverage gaps.} Within the 312-method corpus, state-modification methods exhibit the lowest proportion of strong specifications (58.7\%), reflecting the difficulty of capturing complete frame conditions for methods that mutate object state under conditional logic. The \emph{other} category---event handlers, I/O operations, and callbacks---contains zero methods in this corpus and is therefore not represented in the controlled results; experience on broader Java code suggests this category is the hardest to specify heuristically, and the corresponding generalisability concern is discussed in Section~\ref{sec:threats}. Where loop heuristics produce only a weak invariant (14.7\% of looping methods), test effectiveness is reduced by approximately 15~pp. Extensions to \texttt{while}-loop and \texttt{for-each} analysis (counter detection, accumulator tracking, and variable-relationship inference) together with the standard library specification database mitigate these gaps in part; loops involving nested data structures or non-linear counters remain difficult.

\paragraph{What Z3 cannot validate.} Of the 985 clauses JML-Inferrer emits, 7.1\% are not discharged by OpenJML's Extended Static Checker (Table~\ref{tab:discharge}; Figure~\ref{fig:discharge}). The non-discharged clauses are not arbitrary failures of the inferrer---they fall into three categorically identifiable patterns that exceed the tractability budget of OpenJML and Z3~4.13.4 at the strictness configuration of Section~\ref{subsec:validation}. (i) \emph{Recursive multiplicative growth.} Z3's linear-arithmetic core cannot unfold the inductive product needed to bound \texttt{n * recurse(n-1)} (factorial), \texttt{recurse(n-1) + recurse(n-2)} (Fibonacci), or \texttt{LIT * recurse(n-1)} (power) without an algebraic step the solver does not perform. JML-Inferrer emits a literal parameter bound for these shapes (\texttt{n <= 12} for factorial, \texttt{n <= 46} for Fibonacci, \texttt{n <= 30} for power) so the body discharges under a small-input contract, but the unbounded form remains unprovable. (ii) \emph{Quantified invariants over mutable arrays.} The shape \texttt{for (i=0; i<arr.length; i++) arr[i] *= factor} requires Z3's array theory to preserve a \texttt{(\textbackslash forall int k; i $\le$ k $<$ arr.length; arr[k]} bound\texttt{)} across an array \texttt{store}; three formulations (forall carry-over, literal-bounded factor, frame condition) each exceeded the 240\,s per-clause budget. (iii) \emph{OpenJML for-each translation.} In a for-each accumulator (\texttt{for (int val : arr) total += val;}), the loop variable \texttt{val} is bound from \texttt{arr[<iter-index>]} internally, but the relationship \texttt{val == arr[\textbackslash count]} is not exposed to user JML; the body's per-step overflow check therefore cannot reference the precondition's \texttt{arr[k]} element bound. The remainder of the non-discharged clauses is dominated by rich-precondition timeouts on dispatcher and matrix-access methods, where heuristic pruning (rather than addition) is the likely tractable path forward. These limits are properties of the SMT solver and the JML translation at the present strictness, not of the inferrer's heuristics; they delimit what verification technology can currently certify, not what JML-Inferrer can emit.

\paragraph{Specifications underdetermined by code.} A second class of limitation is structural rather than heuristic: some specifications cannot, even in principle, be inferred from the code alone, because the implementation makes a choice that the abstract contract leaves open. The canonical example is sorting in the presence of duplicate keys. A stable merge sort and an unstable quicksort are both correct realisations of \texttt{void sort(int[] arr)} relative to the abstract contract ``the result is a permutation of the input that is non-decreasing'', but they differ in their treatment of equal elements: the merge sort preserves the input order of equals; the quicksort does not. Code-derived inference of the merge sort body will produce a stability postcondition that is true of \emph{this} implementation but not of the API; a caller that depends on the inferred stability clause becomes coupled to the implementation rather than the contract. Analogous cases include hash-table iteration order, exception-message text, floating-point rounding direction on ties, and any place where the standard library says ``the result is one of \dots''. JML-Inferrer makes no attempt to recover the abstract over-approximation in these cases; the discharge oracle (OpenJML against the implementation) cannot detect the over-specification either, because the inferred clause is consistent with the body by construction. Mitigation requires an external source---documentation, a reference specification, or comparison across multiple implementations of the same interface---and is identified as a follow-up direction in Section~\ref{sec:conclusion}.

\paragraph{Language-feature limitations.} The current treatment of object-oriented constructs has known gaps, including conflicts between inherited specifications, the difficulty that polymorphic dispatch poses for postcondition inference, and the absence of concurrency properties from the analysis. Branch-sensitive postcondition inference (handling \texttt{if}/\texttt{else} returns and ternary expressions) improves postcondition completeness for methods with branching return logic; deeply nested or multi-branch conditions may still yield conservative results.

\paragraph{False positives.} Heuristic pattern matching can in principle emit clauses that are syntactically valid yet semantically inconsistent with the implementation; such clauses would mislead developers or cause spurious test failures. The principal defence is the OpenJML discharge step (Section~\ref{subsec:validation}, Table~\ref{tab:discharge}): 92.9\% of inferred clauses are discharged formally, 3.8\% are flagged for review, and the remaining 3.4\% involve JML constructs that the verifier cannot yet handle. Clauses that fail formal verification are routed to the review queue rather than silently retained, catching residual unsoundness in numeric corner cases such as integer-overflow boundaries and floating-point edge values. The standard library specification database is curated manually, which constrains the inputs to interprocedural propagation. Plausible additional mitigations include confidence scoring and human review for safety-critical methods.

A detailed comparison with related specification inference methods is provided in the Supplementary Material (Appendix~E).
